\documentclass[10pt,a4paper]{amsart}
\usepackage{amsmath,amssymb,amsthm,graphicx}
% For commutative diagrams we use Xy-pic
\usepackage[all]{xy}

% We use 2cell for 2-commutative diagrams.
\xyoption{2cell}
% Theorem environments.
%
\theoremstyle{plain}
\newtheorem{theorem}[subsection]{Theorem}
\newtheorem{proposition}[subsection]{Proposition}
\newtheorem{lemma}[subsection]{Lemma}

\theoremstyle{definition}
\newtheorem{definition}[subsection]{Definition}
\newtheorem{example}[subsection]{Example}
\newtheorem{exercise}[subsection]{Exercise}
\newtheorem{situation}[subsection]{Situation}

\theoremstyle{remark}
\newtheorem{remark}[subsection]{Remark}
\newtheorem{remarks}[subsection]{Remarks}

\numberwithin{equation}{subsection}

% Macros
%
\def\lim{\mathop{\mathrm{lim}}\nolimits}
\def\colim{\mathop{\mathrm{colim}}\nolimits}
\def\Spec{\mathop{\mathrm{Spec}}}
\def\Hom{\mathop{\mathrm{Hom}}\nolimits}
\def\Ext{\mathop{\mathrm{Ext}}\nolimits}
\def\SheafHom{\mathop{\mathcal{H}\!\mathit{om}}\nolimits}
\def\SheafExt{\mathop{\mathcal{E}\!\mathit{xt}}\nolimits}
\def\Sch{\mathit{Sch}}
\def\Mor{\mathop{\mathrm{Mor}}\nolimits}
\def\Ob{\mathop{\mathrm{Ob}}\nolimits}
\def\Sh{\mathop{\mathit{Sh}}\nolimits}
\def\NL{\mathop{N\!L}\nolimits}
\def\CH{\mathop{\mathrm{CH}}\nolimits}
\def\proetale{{pro\text{-}\acute{e}tale}}
\def\etale{{\acute{e}tale}}
\def\QCoh{\mathit{QCoh}}
\def\Ker{\mathop{\mathrm{Ker}}}
\def\Im{\mathop{\mathrm{Im}}}
\def\Coker{\mathop{\mathrm{Coker}}}
\def\Coim{\mathop{\mathrm{Coim}}}

% Boxtimes
%
\DeclareMathSymbol{\boxtimes}{\mathbin}{AMSa}{"02}

%
% Macros for moduli stacks/spaces
%
\def\QCohstack{\mathcal{QC}\!\mathit{oh}}
\def\Cohstack{\mathcal{C}\!\mathit{oh}}
\def\Spacesstack{\mathcal{S}\!\mathit{paces}}
\def\Quotfunctor{\mathrm{Quot}}
\def\Hilbfunctor{\mathrm{Hilb}}
\def\Curvesstack{\mathcal{C}\!\mathit{urves}}
\def\Polarizedstack{\mathcal{P}\!\mathit{olarized}}
\def\Complexesstack{\mathcal{C}\!\mathit{omplexes}}
% \Pic is the operator that assigns to X its picard group, usage \Pic(X)
% \Picardstack_{X/B} denotes the Picard stack of X over B
% \Picardfunctor_{X/B} denotes the Picard functor of X over B
\def\Pic{\mathop{\mathrm{Pic}}\nolimits}
\def\Picardstack{\mathcal{P}\!\mathit{ic}}
\def\Picardfunctor{\mathrm{Pic}}
\def\Deformationcategory{\mathcal{D}\!\mathit{ef}}
\newenvironment{reference}{}{}
\newenvironment{slogan}{}{}
\newenvironment{history}{}{}
\usepackage{cleveref}
\crefname{originallemma}{original lemma}{original lemmas}
\Crefname{originallemma}{Original lemma}{Original lemmas}
\makeatletter
\newlabel{divisors-lemma-characterize-effective-Cartier-divisor}{{01WS}{source}{}{sourcepointer-1}{}}
\newlabel{divisors-lemma-characterize-effective-Cartier-divisor@cref}{{[originallemma][0][]01WS}{[1][0][]source}{}{}{}}
\makeatother
\crefname{originallemma}{original lemma}{original lemmas}
\Crefname{originallemma}{Original lemma}{Original lemmas}
\makeatletter
\newlabel{divisors-lemma-integral-effective-Cartier-divisor-dvr}{{0AG9}{source}{}{sourcepointer-2}{}}
\newlabel{divisors-lemma-integral-effective-Cartier-divisor-dvr@cref}{{[originallemma][0][]0AG9}{[1][0][]source}{}{}{}}
\makeatother
\crefname{originallemma}{original lemma}{original lemmas}
\Crefname{originallemma}{Original lemma}{Original lemmas}
\makeatletter
\newlabel{divisors-lemma-effective-Cartier-divisor-is-a-sum}{{0BCP}{source}{}{sourcepointer-3}{}}
\newlabel{divisors-lemma-effective-Cartier-divisor-is-a-sum@cref}{{[originallemma][0][]0BCP}{[1][0][]source}{}{}{}}
\makeatother
\crefname{originaldefinition}{original definition}{original definitions}
\Crefname{originaldefinition}{Original definition}{Original definitions}
\makeatletter
\newlabel{chow-definition-cycle-associated-to-closed-subscheme}{{02QU}{source}{}{sourcepointer-4}{}}
\newlabel{chow-definition-cycle-associated-to-closed-subscheme@cref}{{[originaldefinition][0][]02QU}{[1][0][]source}{}{}{}}
\makeatother
\crefname{originallemma}{original lemma}{original lemmas}
\Crefname{originallemma}{Original lemma}{Original lemmas}
\makeatletter
\newlabel{more-algebra-lemma-dedekind-torsion-free-flat}{{0AUW}{source}{}{sourcepointer-5}{}}
\newlabel{more-algebra-lemma-dedekind-torsion-free-flat@cref}{{[originallemma][0][]0AUW}{[1][0][]source}{}{}{}}
\makeatother
\title{The special fibre of a regular model is an effective Cartier divisor cut out by a principal ideal, and so is each of its components}
\author{The Stacks Project Authors}
\date{stacks-project:models.tex:local-snapshot-20261012; GNU FDL 1.2 or later}
\begin{document}
\maketitle
\noindent\textit{Editorial note.} The counted unit is the Stacks Project Lemma lemma-divisor-special-fiber at source lines 4547-4558 together with its proof at source lines 4560-4604, a straight adjacent pairing in the chapter file. The chapter is the upstream file models.tex; the pool directory carries the longer name stacks-models and is not used as the chapter name anywhere in this card.\par\medskip
\noindent\textit{Editorial note.} This lemma is the one that five earlier cards of this lane already point at by its tag 0C5Z, notably the cards on the Picard group of a regular model, on the multiple-fibre normal bundle and on the intersection pairing; those pointers and this counted unit are now consistent.\par\medskip
\noindent\textit{Editorial note.} This card has no context span: the statement and proof are self-contained apart from the five pointers, and the notions they use are recorded in the standard prerequisites.\par\medskip
\noindent\textit{Editorial note.} The retained proof quotes five results of the Divisors, Chow and More on Algebra chapters; all five are declared as non-load-bearing navigation pointers with their stable Stacks tags looked up in the shared tag index, and none of their statements is reproduced inside the span.\par\medskip
\noindent\textit{Editorial note.} The retained proof span starts at line 4560, the line of the opening marker read from the source, and contains exactly one proof environment.\par\medskip
\noindent\textit{Editorial note.} No citation occurs in the retained spans, so no bibliography entry is supplied.\par\medskip
\noindent\textit{Allowed standard prerequisites (editorial).}\begin{itemize}
\item Effective Cartier divisors on a regular surface over a discrete valuation ring. A regular model of a smooth curve over the fraction field of a discrete valuation ring is a regular scheme of dimension two, flat and proper over the ring, with a special fibre that is a divisor supported on finitely many irreducible curves; an effective Cartier divisor on such a scheme is cut out locally by a single nonzerodivisor, and a divisor supported on one irreducible component is Cartier when the local ring at its generic point is a discrete valuation ring. A torsion-free module over a Dedekind domain is flat, and a scheme-theoretic cycle may be associated to a closed subscheme of the expected codimension. The author's own notions and the results of the Divisors, Chow and More on Algebra chapters that the proof quotes are the allowed inputs.. The counted lemma is the divisorial description of the special fibre that the chapter uses everywhere afterwards, for instance when it attaches an invertible sheaf to the special fibre and when it speaks of degrees on the components.
\end{itemize}\medskip
\ifdefined\hypertarget\hypertarget{sourcepointer-1}{}\fi
\noindent\textit{Source reference.} Original lemma 01WS: focused-seeds/AGP/sources/stacks-snapshot-20261011/tags.txt, source lines 2493--2493. Allowed input of the Divisors chapter quoted by the retained proof for the characterisation of an effective Cartier divisor, which is what the first two assertions verify; the pointer is registered so the reference resolves and the cited statement is not reproduced inside the span. The number is the stable Stacks tag of the target, looked up in the shared tag index.\par\medskip
\ifdefined\hypertarget\hypertarget{sourcepointer-2}{}\fi
\noindent\textit{Source reference.} Original lemma 0AG9: focused-seeds/AGP/sources/stacks-snapshot-20261011/tags.txt, source lines 13187--13187. Allowed input of the Divisors chapter quoted by the retained proof to recognise a component as Cartier by looking at the local ring at its generic point; the pointer is registered so the reference resolves and the cited statement is not reproduced inside the span. The number is the stable Stacks tag of the target, looked up in the shared tag index.\par\medskip
\ifdefined\hypertarget\hypertarget{sourcepointer-3}{}\fi
\noindent\textit{Source reference.} Original lemma 0BCP: focused-seeds/AGP/sources/stacks-snapshot-20261011/tags.txt, source lines 14287--14287. Allowed input of the Divisors chapter quoted by the retained proof for the description of a sum of effective Cartier divisors, which supplies the displayed expansion of the special fibre; the pointer is registered so the reference resolves and the cited statement is not reproduced inside the span. The number is the stable Stacks tag of the target, looked up in the shared tag index.\par\medskip
\ifdefined\hypertarget\hypertarget{sourcepointer-4}{}\fi
\noindent\textit{Source reference.} Original definition 02QU: focused-seeds/AGP/sources/stacks-snapshot-20261011/tags.txt, source lines 3601--3601. Allowed input of the Chow Homology chapter quoted by the retained proof for the cycle associated to a closed subscheme, used when the multiplicities of the components are read off; the pointer is registered so the reference resolves and the cited definition is not reproduced inside the span. The number is the stable Stacks tag of the target, looked up in the shared tag index.\par\medskip
\ifdefined\hypertarget\hypertarget{sourcepointer-5}{}\fi
\noindent\textit{Source reference.} Original lemma 0AUW: focused-seeds/AGP/sources/stacks-snapshot-20261011/tags.txt, source lines 13664--13664. Allowed input of the More on Algebra chapter quoted by the retained proof for the flatness of a torsion-free module over a Dedekind domain, which underlies the uniformity argument over the base; the pointer is registered so the reference resolves and the cited statement is not reproduced inside the span. The number is the stable Stacks tag of the target, looked up in the shared tag index.\par\medskip
\begin{lemma}
\label{lemma-divisor-special-fiber}
Let $X$ be a regular model of a smooth curve $C$ over $K$.
\begin{enumerate}
\item the special fibre $X_k$ is an effective Cartier divisor on $X$,
\item each irreducible component $C_i$ of $X_k$ is an effective
Cartier divisor on $X$,
\item $X_k = \sum m_i C_i$ (sum of effective Cartier divisors)
where $m_i$ is the multiplicity of $C_i$ in $X_k$,
\item $\mathcal{O}_X(X_k) \cong \mathcal{O}_X$.
\end{enumerate}
\end{lemma}

\begin{proof}
Recall that $R$ is a discrete valuation ring with uniformizer $\pi$
and residue field $k = R/(\pi)$. Because $X \to \Spec(R)$ is flat,
the element $\pi$ is a nonzerodivisor affine locally on $X$
(see More on Algebra, Lemma
\ref{more-algebra-lemma-dedekind-torsion-free-flat}). Thus
if $U = \Spec(A) \subset X$ is an affine open, then
$$
X_k \cap U = U_k = \Spec(A \otimes_R k) = \Spec(A/\pi A)
$$
and $\pi$ is a nonzerodivisor in $A$.
Hence $X_k = V(\pi)$ is an effective Cartier divisor by
Divisors, Lemma \ref{divisors-lemma-characterize-effective-Cartier-divisor}.
Hence (1) is true.

\medskip\noindent
The discussion above shows that the pair $(\mathcal{O}_X(X_k), 1)$
is isomorphic to the pair $(\mathcal{O}_X, \pi)$ which proves (4).

\medskip\noindent
By Divisors, Lemma \ref{divisors-lemma-effective-Cartier-divisor-is-a-sum}
there exist pairwise distinct integral effective Cartier divisors
$D_i \subset X$ and integers $a_i \geq 0$ such that $X_k = \sum a_i D_i$.
We can throw out those divisors $D_i$ such that $a_i = 0$. Then it is
clear (from the definition of addition of effective Cartier
divisors) that $X_k = \bigcup D_i$ set theoretically. Thus $C_i = D_i$
are the irreducible components of $X_k$ which proves (2).
Let $\xi_i$ be the generic point of $C_i$.
Then $\mathcal{O}_{X, \xi_i}$ is a discrete valuation ring
(Divisors, Lemma \ref{divisors-lemma-integral-effective-Cartier-divisor-dvr}).
The uniformizer $\pi_i \in \mathcal{O}_{X, \xi_i}$ is a local equation
for $C_i$ and the image of $\pi$ is a local equation for $X_k$.
Since $X_k = \sum a_i C_i$ we see that $\pi$ and $\pi_i^{a_i}$
generate the same ideal in $\mathcal{O}_{X, \xi_i}$.
On the other hand, the multiplicity of $C_i$ in $X_k$ is
$$
m_i = \text{length}_{\mathcal{O}_{C_i, \xi_i}} \mathcal{O}_{X_k, \xi_i} =
\text{length}_{\mathcal{O}_{C_i, \xi_i}} \mathcal{O}_{X, \xi_i}/(\pi) =
\text{length}_{\mathcal{O}_{C_i, \xi_i}} \mathcal{O}_{X, \xi_i}/(\pi_i^{a_i}) =
a_i
$$
See Chow Homology, Definition
\ref{chow-definition-cycle-associated-to-closed-subscheme}.
Thus $a_i = m_i$ and (3) is proved.
\end{proof}

\end{document}
